% Complete Phase II proofs. Parent manuscript supplies theorem environments.
\section{What an activation-only controller cannot observe}
\label{sec:activation-observation}

We now study actual neural parameters. This model has no discrete
representation-cardinality constraint. Its purpose is to separate an
information summary from the information needed to realize a structural
operation. The observation restriction below is explicit: a physical unit
may also have access to presynaptic inputs, in which case the negative
result does not apply.

Let $X$ be uniform on $\{e_1,-e_1,e_2,-e_2\}\subset\mathbb R^2$.
For an unknown task $s\in\{-1,+1\}$ and $0<\eta<1/2$, set
\[
 q_s(\pm e_1)=\mathbb P_s(Y=1\mid X=\pm e_1)=\tfrac12+s\eta,
 \qquad q_s(\pm e_2)=\tfrac12-s\eta.
\]
Every source-label pair has positive probability. Write
$\phi(t)=\tanh t$ and use binary log loss
$\ell(f,y)=\log(1+e^f)-yf$. A one-hidden-unit network has logit
$f(x)=c+a\phi(b+w^\top x)$. Fix $b=b_0>0$, $a=1$, $w=0$, and
$c=-\phi(b_0)$, so that $f_0=0$ and the predicted probability is $1/2$.
At this point the activation is $A=\phi(b_0)$ and the output residual is
$R=1/2-Y$. All logs below are natural.

\begin{proposition}[A parameter local minimum with unrepresented task information]
\label{prop:axis-local-min}
If
\begin{equation}
 0<\eta<\frac{\phi'(b_0)^2}{4|\phi''(b_0)|},
 \label{eq:axis-stability}
\end{equation}
the stated one-unit parameter vector is a local minimum of population
log loss on either task. It lies on a manifold of equivalent constant
predictors. Nevertheless
\[
 I_s(X;Y)=\log2-h(\tfrac12+\eta)>0,
 \qquad I_s(A;Y)=0.
\]
The entire laws of $(A,Y,R)$ are identical for the two tasks.
\end{proposition}
\begin{proof}
Introduce the smooth invertible coordinate $r=c+a\phi(b)$, retaining
$a,b,w$. Then
$f(x)=r+a\{\phi(b+w^\top x)-\phi(b)\}$.
For every $a,b$ near $1,b_0$, the risk as a function of $(r,w)$ is
stationary at $(0,0)$: its $r$ derivative is
$\mathbb E_s(1/2-Y)=0$, and its $w$ derivative is
$a\phi'(b)\mathbb E_s[(1/2-Y)X]=0$ by sign symmetry of each axis.
At this point its Hessian in the $(r,w)$ coordinates is block diagonal,
with $rr$ entry $1/4$ and $ww$ block
\begin{equation}
 \frac{a^2\phi'(b)^2}{8}I_2
 +a\phi''(b)
 \begin{pmatrix}-s\eta/2&0\\0&s\eta/2\end{pmatrix}.
 \label{eq:axis-hessian}
\end{equation}
Indeed $\ell_f=1/2-Y$, $\ell_{ff}=1/4$ at zero,
$\mathbb E XX^\top=I_2/2$, and the mixed $rw$ entries vanish because
$\mathbb EX=0$. Under \eqref{eq:axis-stability} this Hessian is
positive definite at $(a,b)=(1,b_0)$. By continuity there is a compact
neighborhood of $(1,b_0,0,0)$ and $m>0$ on which the $(r,w)$ Hessian is
bounded below by $mI_3$. Taylor's integral formula along the line from
$(0,0)$ to $(r,w)$, with $a,b$ fixed, therefore gives
\[
 L_s(r,w;a,b)-\log2\ge \frac m2(r^2+\|w\|^2)
\]
throughout a smaller neighborhood. At $r=w=0$ the risk remains $\log2$
for all nearby $a,b$. This proves a local minimum in all original
parameters, allowing its neutral directions.

The marginal law of $Y$ is Bernoulli$(1/2)$ and every input has
conditional entropy $h(1/2+\eta)$; the information formula follows.
Since $A$ is constant, and $R$ is a function of $Y$, the final assertions
follow as well.
\end{proof}

Replace the hidden unit by two units of output weight $1/2$, biases
$b_0$, and incoming weights $+\epsilon v,-\epsilon v$, retaining the
intercept $-\phi(b_0)$, where $\|v\|=1$. The resulting logit is
\begin{equation}
 f_{\epsilon,v}(x)=\frac{\phi(b_0+\epsilon v^\top x)
 +\phi(b_0-\epsilon v^\top x)}2-\phi(b_0).
 \label{eq:axis-split}
\end{equation}
At $\epsilon=0$ this is a function-preserving duplication; at positive
$\epsilon$ it is the specified structural perturbation. This familiar
symmetric split is the operation studied in splitting steepest descent;
the score derived below is not claimed as new.

For $t\ge0$ put
\begin{equation}
 d(t)=\phi(b_0)-\frac{\phi(b_0+t)+\phi(b_0-t)}2
 =\frac{z(1-z^2)\tanh^2t}{1-z^2\tanh^2t},\qquad z=\tanh b_0.
 \label{eq:axis-d}
\end{equation}
Thus $d(0)=0$ and $0<d(t)<z$ for $t>0$.

\begin{theorem}[Exact observation-channel obstruction and attainable descent]
\label{thm:axis-observation}
Let $\Delta_s(\epsilon,v)=L_s(f_{\epsilon,v})-\log2$, and put
$d_j=d(|\epsilon v_j|)$. Then
\begin{equation}
 \Delta_s(\epsilon,v)
 =\frac12\left\{\log\cosh(d_1/2)+\log\cosh(d_2/2)\right\}
   +\frac{s\eta}{2}(d_1-d_2).
 \label{eq:axis-risk}
\end{equation}
Consequently:
\begin{enumerate}
\item Any possibly randomized controller whose entire task-dependent
observation consists of the population law of $(A,Y,R)$, or any number
of independent samples from that law, satisfies
\[
 \max_{s\in\{-1,+1\}}\mathbb E_s
       \Delta_s(\epsilon,v)\ge0.
\]
Here it may select both $\epsilon\ge0$ and the unit vector $v$ from
these observations. Expectations include its observations and randomness;
risk is evaluated on a fresh population input. The controller may also
decline to split, which gives zero change.
\item A controller with the input-residual second moment
$M_s=\mathbb E_s[(1/2-Y)XX^\top]$ can choose $v=e_2$ for $s=+1$
and $v=e_1$ for $s=-1$. For either task its exact gain is
\begin{equation}
 G(\epsilon,\eta)
 =\frac{\eta d(\epsilon)}2
  -\frac12\log\cosh\{d(\epsilon)/2\}
 \ge\frac{\eta d(\epsilon)}2-\frac{d(\epsilon)^2}{16}>0
 \label{eq:axis-gain}
\end{equation}
whenever $0<d(\epsilon)<8\eta$.
\end{enumerate}
\end{theorem}
\begin{proof}
The addition formulas for $\tanh(b_0\pm t)$ give
\eqref{eq:axis-d}; its denominator is positive. At either point on
axis $j$, the split logit is $-d_j$. For any posterior $q$,
\[
 \mathbb E[\ell(-d,Y)\mid q]-\log2
 =\log\cosh(d/2)+(q-1/2)d.
\]
Averaging the two axes proves \eqref{eq:axis-risk}.

The two tasks give identical laws of all permitted controller
observations. Consequently its joint distribution of $(\epsilon,v)$
is the same under both tasks, including when selected from a random
sample. Pointwise,
\[
 \Delta_+(\epsilon,v)+\Delta_-(\epsilon,v)
 =\log\cosh(d_1/2)+\log\cosh(d_2/2)\ge0.
\]
Taking the common expectation proves the first part. Independent
controller randomization cannot change this calculation.

The matrix is $M_s=\operatorname{diag}(-s\eta/2,s\eta/2)$, so it
identifies the displayed axis choice. Substituting it in
\eqref{eq:axis-risk} gives the gain. To verify the inequality used
in \eqref{eq:axis-gain}, for $u\ge0$ one has $\tanh u\le u$, because
the derivative of $u-\tanh u$ is $\tanh^2u\ge0$ and the value at zero
is zero. Integrating gives $\log\cosh u\le u^2/2$; evenness handles
$u<0$. The strict positivity follows from $d<8\eta$.
\end{proof}

For orientation, Taylor expansion of \eqref{eq:axis-split} yields
\[
 \Delta_s(\epsilon,v)
 =\frac{\epsilon^2}{2}v^\top S_s v+O(\epsilon^4),
 \qquad S_s=\phi''(b_0)M_s.
\]
This is the established splitting matrix specialized to log loss.
Theorem~\ref{thm:axis-observation} adds an exact finite-displacement
statement about which observation channels can select its useful
direction. It does not require naming $M_s$ an entropy estimator or a
unique-information decomposition. In particular, even the scalar
$I_s(X;Y)$ is identical across the tasks and does not specify an axis.

\begin{theorem}[Finite observations: sufficient and necessary sample order]
\label{thm:axis-samples}
Suppose the controller receives $n$ independent full pairs $(X_i,Y_i)$.
Define
\[
 U_i=(1/2-Y_i)(X_{i2}^2-X_{i1}^2),\qquad
 \widehat s=\operatorname{sign}\Big(\sum_{i=1}^n U_i\Big),
\]
with any fixed rule for a zero sum. Select the axis prescribed in
Theorem~\ref{thm:axis-observation} for $\widehat s$, with fixed
$\epsilon>0$ satisfying $d(\epsilon)<8\eta$. On either task,
\begin{equation}
 \mathbb P_s\{\Delta_s=-G(\epsilon,\eta)\}
 \ge1-\exp(-2n\eta^2).
 \label{eq:axis-sample-upper}
\end{equation}
Conversely, for any randomized procedure selecting a symmetric split
from $n$ full observations, if its probability of strict population
descent is at least $1-\delta$ on each task, where $0<\delta<1/4$,
then
\begin{equation}
 n\ge\frac{\log(1/(4\delta))}{-\log(1-4\eta^2)}.
 \label{eq:axis-sample-lower}
\end{equation}
For $0<\eta\le1/4$, the right hand side is at least
$3\log(1/(4\delta))/(16\eta^2)$.
\end{theorem}
\begin{proof}
For task $s$, $U_i$ takes values $\pm1/2$ and has mean $s\eta$.
Hoeffding's inequality for independent variables with range length one
therefore gives
$\mathbb P_s(s\sum_i U_i\le0)\le e^{-2n\eta^2}$.
For completeness its exponential bound follows by applying Markov's
inequality to $e^{-t\sum_i(sU_i-\eta)}$, using the bounded-variable
moment generating function inequality
$\mathbb Ee^{t(V-\mathbb EV)}\le e^{t^2/8}$ for an interval of length
one, and minimizing $e^{-tn\eta+nt^2/8}$ at $t=4\eta$.
The moment generating function inequality follows from convexity by
bounding $e^{tV}$ by its endpoint chord and maximizing the centered
two-point expression (or by twice integrating that expression's
second logarithmic derivative, which is a Bernoulli variance at most
$1/4$). Correct identification gives exactly the claimed gain.

Any selected split with $\Delta_s<0$ must, by \eqref{eq:axis-risk},
satisfy $s(d_2-d_1)>0$. Therefore a procedure that achieves strict
descent with probability at least $1-\delta$ yields a test for $s$
with each error probability at most $\delta$: report $+1$ when
$d_2>d_1$ and $-1$ otherwise. Extra independent randomization cannot
improve the optimal testing error.

Write $P_+,P_-$ for the two laws of a full pair. Their Hellinger
affinity is
\[
 \rho=\sum_{x,y}\sqrt{P_+(x,y)P_-(x,y)}
     =\sqrt{1-4\eta^2}.
\]
The affinity of product laws is $\rho^n$. For arbitrary finite laws
$P,Q$ with affinity $a$, Cauchy--Schwarz gives
\[
 \|P-Q\|_{\rm TV}
 =\frac12\sum_z|\sqrt{P(z)}-\sqrt{Q(z)}|
                 (\sqrt{P(z)}+\sqrt{Q(z)})
 \le\sqrt{1-a^2}.
\]
The sum of type-I and type-II errors of any test is at least
$1-\|P-Q\|_{\rm TV}$: this follows directly by writing its acceptance
probability as a function in $[0,1]$ and bounding the difference of
expectations by total variation. Its equal-prior average error is
thus at least
\[
 \frac{1-\sqrt{1-\rho^{2n}}}{2}\ge\frac{\rho^{2n}}4
 =\frac{(1-4\eta^2)^n}4,
\]
where the inequality uses $1-\sqrt{1-t}=t/(1+\sqrt{1-t})\ge t/2$.
Both errors at most $\delta$ imply \eqref{eq:axis-sample-lower}.
Finally $-\log(1-u)=\int_0^u(1-t)^{-1}dt\le u/(1-u)$ and
$u=4\eta^2\le1/4$ imply
$-\log(1-4\eta^2)\le16\eta^2/3$.
\end{proof}

The lower bound applies even to a controller given full inputs; the
activation-only controller receives no task-identifying information at
any sample size. The upper bound uses ordinary independent sampling,
not an oracle population covariance. Neither statement predicts the
outcome of unrestricted training after the operation. Such training
observes input coordinates and can learn distinctions unavailable to
the restricted structural controller.

